% !TeX root = ../bnuthesis-example.tex

\begin{bnuindex}

\section*{索引说明}
本学位论文编制了创新索引、主题索引、专有名称索引。

\section*{创新索引}
\addcontentsline{toc}{section}{创新索引}
创新索引以学位论文中描述创新内容的词语为目标编制，款目采用\textbf{加粗}格式表示，按标目在论文中首次出现的顺序集中排列在索引的最前面。

\vspace{1em}
% 创新索引内容
\textbf{蓝莓发酵酒最佳发酵条件（工艺参数）}\dotfill\pageref{idx:innov:ferment}

\textit{最佳原材料为栽培蓝莓，糖种类为蜂蜜，酵母菌种为 Fermivin。最佳发酵温度为 15 ℃，初始糖度为 26\%，酵母添加量为 0.3\%，硫代硫酸钠添加量为 100 ppm。}

\textbf{蓝莓发酵酒抗氧化特性}\dotfill\pageref{idx:innov:antioxi}

\textbf{蓝莓发酵酒降血脂作用}\dotfill\pageref{idx:innov:lipid}

\textbf{蓝莓发酵酒降细胞毒性作用}\dotfill\pageref{idx:innov:lipid}, \pageref{idx:innov:antioxi} 


\section*{主题索引}
\addcontentsline{toc}{section}{主题索引}
主题索引以学位论文中重点论述、具有检索价值的重要主题词为标目编制，按拼音排序。

\vspace{1em}
% 主题索引内容
\indexletter{C}
\indexentry{超氧化物歧化酶(SOD)}{idx:theme:superoxy} 
\indexentry{超氧阴离子}{idx:theme:superoxy}
\indexentry{初始糖度}{idx:theme:sugar1} 

\indexletter{G}
\indexentry{感官评价}{idx:innov:ferment} 
\indexentry{果胶含量测定}{idx:theme:pectin}


\section*{专有名称索引}
\addcontentsline{toc}{section}{专有名称索引}
专有名称索引以机构名、人名、地名、文献名等专有名称为标目编制，按拼音排序。

\vspace{1em}
% 专有名称索引内容
\indexletter{A}
\indexentry{安徽广播电视台}{idx:name:anhuitv}
\indexentry{《安徽农业科学》}{idx:name:anhuiagr} 

\indexletter{G}
\indexentry{郭嘉欣}{idx:name:guojiaxin}
\indexentry{《国家安全太空战略》}{idx:innov:lipid}

\end{bnuindex}